% Vista científica exacta materializada desde una rama condicional.
% Fuente: source/demostraciones_primarias/27_04_particula_estadistica.tex; rama: HMTIncludeParticleCore.
% No modifica el contenido matemático de la rama de origen.
\chapter{Secciones persistentes y propagación finita}
\label{chap:particula}

La noción matemática de partícula se introduce después de la estructura de
frontera y antes de especializar su dinámica. Su soporte primario es una
extensión superviviente del sistema de estados TPK\@. Una trayectoria
orientada es la sombra secuencial de esa extensión; sobre ella, los grados
internos forman un fibrado discreto. El propagador, la suma de caminos y las
representaciones de Clifford se incorporan después de construir la recta de
acción. Esta jerarquía impide que una partícula se defina mediante la elección
externa de una fila o de una ruta y separa el objeto matemático de su posterior
identificación con una especie física.

\section{Extensión superviviente y proyección observable}

Sea \((\mathcal S_m,\rho_{m+1,m})\) el sistema inverso de estados que
superan los cierres hasta profundidad \(m\). Un estado global es una familia
compatible
\[
 x=(x_m)_m\in\Omega_{\HMT}^{\rm surv}
 :=\varprojlim_m\mathcal S_m.
\]
Una carta secuencial de profundidad \(R\) produce una ruta
\[
 \gamma_x^{[R]}=\operatorname{sh}_R(x)
 =(c_0\longrightarrow\cdots\longrightarrow c_R).
\]
El símbolo \(\operatorname{sh}\) denota una lectura de sombra: diferentes
cartas pueden presentar un mismo estado global, mientras conserven hoja,
orientación, borde y registro. En el
\cref{chap:extension-ruta-caracter} se construye el registro enriquecido que
hace descender el carácter de acción desde estas presentaciones a la
extensión.

\section{Partícula como sección compatible de un fibrado discreto}

Sea \(\Gamma_{\rm TPK}\) el grafo cuyos vértices son estados completos del
núcleo de fase topológica y cuyas aristas son transformaciones admisibles.
Para \(x\in\Omega_{\HMT}^{\rm surv}\) y su trayectoria observada
\[
 \gamma_x=(c_0\longrightarrow c_1\longrightarrow\cdots
 \longrightarrow c_R),
\]
sea \(\operatorname{Typ}(c_t)\) el conjunto finito de tipos aritméticos,
orientaciones y memorias sobre el estado \(c_t\), y sea
\(E_{c_t}=\mathbb C[\operatorname{Typ}(c_t)]\) el espacio vectorial complejo
libre generado por esos tipos. El fibrado discreto restringido a \(\gamma_x\) es
\[
 \mathcal B_{\gamma_x}=\bigsqcup_{t=0}^{R}E_{c_t}
 \longrightarrow\{0,\ldots,R\}.
\]
Cada arista \(c_t\to c_{t+1}\) induce una aplicación lineal de transporte
\(\tau_t:E_{c_t}\to E_{c_{t+1}}\).

\begin{definition}[Partícula matemática]
Una partícula de Mecánica Dimensional es una tupla
\[
 \mathfrak p=(x,\gamma_x,\mathcal B_{\gamma_x},s,n,U),
\]
donde \(x\in\Omega_{\HMT}^{\rm surv}\) es la extensión superviviente y
\(\gamma_x=\operatorname{sh}_R(x)\) su sombra en la carta declarada; además,
\[
 s:\{0,\ldots,R\}\longrightarrow\mathcal B_{\gamma_x},\qquad
 s(t)\in E_{c_t},\qquad s(t+1)=\tau_t(s(t)),
\]
es una sección compatible;
\[
 n:\bigsqcup_{t=0}^{R}\operatorname{Modes}(E_{c_t})
 \longrightarrow\mathbb N
\]
es su función de ocupación, donde
\(\operatorname{Modes}(E_{c_t}):=\operatorname{Typ}(c_t)\) denota la base
distinguida del espacio vectorial libre; y
\[
 U_{t_2,t_1}:E_{c_{t_1}}\longrightarrow E_{c_{t_2}}
\]
es una familia de propagadores lineales que entrelaza los transportes,
\(U_{t,t}=\operatorname{id}_{E_{c_t}}\),
\(U_{t+1,t}=\tau_t\), y satisface
\[
 U_{t_3,t_2}\circ U_{t_2,t_1}=U_{t_3,t_1}
 \qquad(0\leq t_1\leq t_2\leq t_3\leq R).
\]
Cada composición está definida entre las fibras correspondientes.
La identificación de \(\mathfrak p\) con una especie
física requiere, además, una representación del grupo de simetría pertinente
y un mapa desde estados HMT hacia los estados de dicha representación.
\end{definition}

La extensión conserva la genealogía y el registro; la ruta expone su lectura
ordenada; los grados internos pertenecen a las fibras, la ocupación a \(n\) y
la evolución a \(U\). La partícula es la compatibilidad de estos niveles, no
uno de ellos por separado. En particular, la ruta no es escogida por la
partícula: es una proyección de la extensión que la realiza.
